%%%%%%%%%%%%%%%%%%%%%%%%%%%%%%%%%%%%%%%%%%%%%%%%%%%%%%%%%%%%%%%%%%%%%%%%%%%%
%% Section 2: Classification
%%%%%%%%%%%%%%%%%%%%%%%%%%%%%%%%%%%%%%%%%%%%%%%%%%%%%%%%%%%%%%%%%%%%%%%%%%%%

\begin{frame}{Ordinary vs Partial Differential Equations}
\textbf{Ordinary Differential Equations (ODE)}

Unknown function depends on \textbf{one} variable.

Example:
\[
\frac{dy}{dx} = x^2
\]

\textbf{Partial Differential Equations (PDE)}

Unknown function depends on \textbf{several} variables.

Example:
\[
\frac{\partial u}{\partial t} = k \frac{\partial^2 u}{\partial x^2}
\]
\end{frame}

\begin{frame}{Order of a Differential Equation}
The \textbf{order} of a differential equation is the highest derivative appearing in it.

First order:
\[
\frac{dy}{dx} = x
\]

Second order:
\[
\frac{d^2y}{dx^2} + y = 0
\]

Third order:
\[
\frac{d^3y}{dx^3} + y = 0
\]
\end{frame}

\begin{frame}{Linear vs Nonlinear Equations}
\textbf{Linear differential equation}
\[
a_n(x)y^{(n)} + a_{n-1}(x)y^{(n-1)} + \dots + a_1(x)y' + a_0(x)y = f(x)
\]

Example:
\[
y'' + 3y' + 2y = 0
\]

\textbf{Nonlinear examples}
\[
y' = y^2 \qquad\qquad y'' + y^3 = 0
\]

\vspace{0.3cm}
\textbf{Why does linearity matter?}
\begin{itemize}
    \item \textbf{Superposition}: If $y_1$ and $y_2$ solve a linear homogeneous equation, so does $c_1 y_1 + c_2 y_2$. You can build complex solutions from simple pieces.
    \item \textbf{Structure}: Linear equations have predictable solution spaces (vector spaces of dimension $n$ for an $n$th-order equation).
    \item \textbf{No chaos}: First-order autonomous linear ODEs cannot exhibit chaotic behaviour --- nonlinearity is \emph{required} for chaos.
\end{itemize}
\end{frame}

\begin{frame}{Homogeneous vs Nonhomogeneous}
\textbf{Linear Homogeneous}
\[
y'' + y = 0
\]

\textbf{Linear Nonhomogeneous}
\[
y'' + y = \sin x
\]

\textbf{Why this matters:}

\begin{itemize}
    \item A homogeneous equation describes a \textbf{free system} (no external driving force).
    \item A nonhomogeneous equation describes a \textbf{driven system} (with an external input $f(x)$).
    \item \textbf{Superposition principle}: the general solution of $Ly = f$ is
    \[
    y = y_h + y_p
    \]
    where $y_h$ solves the homogeneous part and $y_p$ is any particular solution.
\end{itemize}
\end{frame}

\begin{frame}{Autonomous vs Nonautonomous}
\textbf{Autonomous (no explicit $t$)}
\[
\frac{dx}{dt} = f(x)
\]

Example:
\[
\frac{dx}{dt} = x(1-x)
\]

\textbf{Non-autonomous}
\[
\frac{dx}{dt} = f(x,t)
\]

Example:
\[
\frac{dx}{dt} = x + \sin t
\]

\vspace{0.3cm}
\textit{Autonomous systems are time-translation invariant: if $x(t)$ is a solution, so is $x(t - t_0)$ for any shift $t_0$.}
\end{frame}

\begin{frame}{Systems of Differential Equations}
Example system:

\[
\begin{cases}
\dot{x} = y \\
\dot{y} = -x
\end{cases}
\]

Vector form:
\[
\frac{d}{dt}
\begin{pmatrix}
 x \\
 y
\end{pmatrix}
=
\begin{pmatrix}
 y \\
 -x
\end{pmatrix}
\]

\vspace{0.3cm}
\textit{Any single higher-order ODE can be rewritten as a system of first-order ODEs. For example, $y'' + y = 0$ becomes $\dot{y}_1 = y_2,\; \dot{y}_2 = -y_1$.}
\end{frame}


\begin{frame}{Initial vs Boundary Value Problems}
\textbf{Initial Value Problem (IVP)}
\begin{itemize}
    \item \textbf{The Math:} $y' = y, \quad y(0) = 1$
    \item \textbf{The Physics:} Evolution over time. You know the starting state at $t=0$, and the equation drives it forward. (e.g., radioactive decay, launching a rocket).
    \item \textbf{Solvability:} \textit{Unambiguous.} As long as the system is well-behaved, theorems guarantee exactly one unique path forward into the future.
\end{itemize}

\vspace{0.4cm}

\textbf{Boundary Value Problem (BVP)}
\begin{itemize}
    \item \textbf{The Math:} $y'' + y = 0, \quad y(0)=0, \quad y(\pi)=0$
    \item \textbf{The Physics:} Spatial constraints. You know the conditions at the physical edges of a domain. (e.g., a guitar string pinned at both ends, or the temperature at two ends of a metal rod).
    \item \textbf{Solvability:} \textit{Tricky!} BVPs can have one unique solution, no solution at all, or infinitely many. (The example above has infinitely many: $y = c \sin(x)$).
\end{itemize}
\end{frame}
